\documentclass[UTF8,zihao=-4,openany,fontset=none]{ctexrep}

\usepackage[a4paper,top=2.5cm,bottom=2.5cm,left=2.6cm,right=2.6cm,headheight=15pt]{geometry}
\usepackage{fontspec}
\usepackage{xeCJK}
\setmainfont{Times New Roman}
\setsansfont{Arial}
\setmonofont{Consolas}
\setCJKmainfont[AutoFakeBold=2.5,AutoFakeSlant=0.15]{SimSun}
\setCJKsansfont[AutoFakeBold=2.5]{SimSun}
\setCJKmonofont{SimSun}
\ctexset{
  chapter={format=\LARGE\bfseries,number=\chinese{chapter},name={第,章}},
  section={format=\Large\bfseries},
  subsection={format=\large\bfseries}
}

\usepackage{amsmath,amssymb,bm}
\usepackage{booktabs,longtable,tabularx,array,multirow}
\usepackage{graphicx}
\usepackage{caption,subcaption,float,placeins}
\usepackage{xcolor}
\usepackage{enumitem}
\usepackage{listings}
\usepackage{fancyhdr,lastpage}
\usepackage{tikz}
\usetikzlibrary{arrows.meta,positioning,calc,shapes.geometric,fit,backgrounds}
\usepackage[hidelinks,unicode]{hyperref}
\usepackage[nameinlink,noabbrev]{cleveref}

\definecolor{ReportBlue}{HTML}{2457A6}
\definecolor{ReportOrange}{HTML}{D9782D}
\definecolor{ReportGold}{HTML}{B28A20}
\definecolor{ReportPink}{HTML}{B44D73}
\definecolor{ReportInk}{HTML}{252A34}
\definecolor{ReportGrey}{HTML}{6B7280}
\definecolor{ReportLight}{HTML}{EEF3FA}
\definecolor{ReportWarm}{HTML}{FFF7ED}

\graphicspath{{figures/}}
\captionsetup{font=small,labelfont=bf,labelsep=quad}
\setlength{\parindent}{2em}
\setlength{\parskip}{0.35em}
\setlength{\emergencystretch}{3em}
\renewcommand{\arraystretch}{1.28}
\pagestyle{fancy}
\fancyhf{}
\fancyhead[L]{\small 小鸡与人眼 650 nm 全视网膜照明仿真实验报告}
\fancyfoot[C]{第 \thepage\ 页}

\lstdefinestyle{workflow}{
  basicstyle=\ttfamily\small,
  backgroundcolor=\color{ReportLight},
  frame=single,
  rulecolor=\color{ReportBlue!45},
  breaklines=true,
  columns=fullflexible,
  showstringspaces=false,
  xleftmargin=1em,
  xrightmargin=1em
}
\lstset{style=workflow}

\newcommand{\D}{\,\mathrm{D}}
\newcommand{\mm}{\,\mathrm{mm}}
\newcommand{\um}{\,\mu\mathrm{m}}
\newcommand{\degree}{\ensuremath{^{\circ}}}
\newcommand{\sourcefile}[1]{\texttt{#1}}
\newcommand{\keyresult}[1]{\textcolor{ReportBlue}{\textbf{#1}}}

% Generated by generate_report_data.py; do not edit manually.
\newcommand{\AdultSixtyDDiameter}{5.988}
\newcommand{\AdultOneTwentyDDiameter}{2.994}
\newcommand{\ChickSixtyDDiameter}{5.952}
\newcommand{\ChickOneTwentyDDiameter}{2.976}
\newcommand{\AdultInfinityAngle}{20.585}
\newcommand{\ChickInfinityAngle}{20.463}
\newcommand{\ZosMaxError}{4.441e-13}
\newcommand{\ZosUnfocusedSpread}{0.82665}
\newcommand{\FocusedCapture}{100.00}
\newcommand{\FocusedUniformity}{0.846}
\newcommand{\DefocusedCapture}{77.13}
\newcommand{\DefocusedUniformity}{0.831}
\newcommand{\HeadlineRows}{60 & 16.7 & 5.952 (超限) & 5.988 (超限) & 5.988 (超限) \\
70 & 14.3 & 5.102 (超限) & 5.133 (超限) & 5.133 (超限) \\
80 & 12.5 & 4.464 (超限) & 4.491 (超限) & 4.491 (超限) \\
90 & 11.1 & 3.968 (超限) & 3.992 (超限) & 3.992 (超限) \\
100 & 10.0 & 3.571 (超限) & 3.593 (超限) & 3.593 (超限) \\
110 & 9.1 & 3.247 (超限) & 3.266 (超限) & 3.266 (超限) \\
120 & 8.3 & 2.976 (超限) & 2.994 (超限) & 2.994 (超限) \\}
\newcommand{\AdultLensRows}{0 & 10.000 & 35.928 & 1013.8 & 可行 \\
-1 & 10.766 & 36.308 & 1035.3 & 可行 \\
-3 & 12.252 & 37.066 & 1079.1 & 可行 \\
-5 & 13.678 & 37.825 & 1123.7 & 可行 \\
-10 & 17.004 & 39.722 & 1239.2 & 可行 \\
-15 & 20.028 & 41.619 & 1360.4 & 超限 \\
-20 & 22.787 & 43.516 & 1487.3 & 超限 \\}
\newcommand{\DefocusRows}{30--45日龄小鸡 & 3.5 & 0.294 & 18.457 & 22.468 \\
6岁儿童 & 5.0 & 0.835 & 17.721 & 23.450 \\
18岁成人 & 5.0 & 0.835 & 17.721 & 23.450 \\}
\newcommand{\AxialRows}{30--45日龄小鸡 & 3.5 & -1.2 & 12.730 & 0.374 \\
6岁儿童 & 5.0 & -0.5 & 1.342 & 0.112 \\
18岁成人 & 5.0 & -0.5 & 1.342 & 0.112 \\}
\newcommand{\ZosRows}{chick\_d10 & 是 & 1.5000 & 1.5000 & 2.220e-13 & 0.000 & 0.00000 \\
child\_d10 & 是 & 3.0000 & 3.0000 & 0.000e+00 & 0.000 & 0.00000 \\
adult\_d10 & 是 & 3.0000 & 3.0000 & 0.000e+00 & 0.000 & 0.00000 \\
adult\_d5\_ext\_m5 & 是 & 3.0000 & 3.0000 & 4.441e-13 & 0.000 & 0.00000 \\
adult\_d10\_ext\_m10 & 是 & 3.0000 & 3.0000 & 0.000e+00 & 0.000 & 0.00000 \\
adult\_d10\_unaccommodated & 否 & 3.0000 & 3.0000 & 0.000e+00 & 292.858 & 0.82665 \\}


\title{小鸡与人眼 650 nm 全视网膜照明\\光学仿真实验报告}
\author{可复现计算模型与 Ansys Zemax OpticStudio ZOS-API 交叉验证}
\date{2026年8月14日\quad 报告版本 1.1.0}

\begin{document}
\pagenumbering{Roman}

\begin{titlepage}
  \centering
  \vspace*{2.2cm}
  {\Huge\bfseries 小鸡与人眼 650 nm 全视网膜照明\par}
  \vspace{0.45cm}
  {\Huge\bfseries 光学仿真实验报告\par}
  \vspace{1.1cm}
  {\Large 独立 ABCD 近轴模型、蒙特卡洛照度评估\par}
  \vspace{0.25cm}
  {\Large 与 Ansys Zemax OpticStudio 24.1 交叉验证\par}
  \vspace{1.5cm}
  \begin{tikzpicture}[x=1cm,y=1cm]
    \draw[ReportGrey,thin] (-6,0)--(6,0);
    \draw[ReportOrange,very thick] (-5,-1.2)--(-5,1.2);
    \node[below] at (-5,-1.35) {圆形面光源};
    \draw[ReportBlue,fill=ReportLight,thick]
      (-1,-1.25) .. controls (-0.45,-0.65) and (-0.45,0.65) .. (-1,1.25)
      .. controls (-1.55,0.65) and (-1.55,-0.65) .. cycle;
    \node[below] at (-1,-1.45) {等效眼透镜};
    \draw[ReportPink,very thick] (4,-1.5) arc[start angle=-90,end angle=90,x radius=0.55,y radius=1.5];
    \node[below] at (4,-1.7) {目标后极部};
    \draw[ReportGold,thick] (-5,1.05)--(-1,0.55)--(4,1.22);
    \draw[ReportBlue,thick] (-5,0)--(-1,0)--(4,0);
    \draw[ReportGold,thick] (-5,-1.05)--(-1,-0.55)--(4,-1.22);
  \end{tikzpicture}
  \vfill
  {\large 技术报告｜中文正文：宋体｜A4\par}
  \vspace{0.3cm}
  {\large 2026年8月14日\par}
  \vspace{1.2cm}
\end{titlepage}

\chapter*{技术摘要}
\addcontentsline{toc}{chapter}{技术摘要}
本研究依据用户提供的《小鸡和人眼光学仿真参数》建立三个旋转对称等效眼模型，研究波长为 $650\,\mathrm{nm}$ 的均匀圆形朗伯面光源在不同物距、瞳孔、离焦、眼轴和外置负镜片条件下覆盖指定后极部区域所需的尺寸。独立模型采用光线状态 $[y,n\theta]^\mathsf{T}$ 与 ABCD 矩阵，在调焦条件 $B=0$ 下求解眼的附加调节力、横向放大率与光源直径；再用每个案例 $600{,}000$ 条确定性随机光线计算相对照度均匀性。最后，ZOS-API 在已安装且 API 许可证有效的 Ansys Zemax OpticStudio 24.1 中自动创建顺序模式 Paraxial 系统并执行批量光线追迹。

主要结论如下：
\begin{itemize}[leftmargin=2.2em]
  \item 物方需求主扫描已改为 $60,70,80,90,100,110,120\D$。无外置镜片时，成人与儿童所需光源直径从 $60\D$ 的 \keyresult{\AdultSixtyDDiameter\,mm} 降至 $120\D$ 的 \keyresult{\AdultOneTwentyDDiameter\,mm}；小鸡分别为 \keyresult{\ChickSixtyDDiameter\,mm} 与 \keyresult{\ChickOneTwentyDDiameter\,mm}。
  \item 三种眼模型在全部 $60$--$120\D$ 请求工况下均超过源文件给定的调节上限，因此上述尺寸是几何覆盖解，而不是生理可调焦解。
  \item 无穷远入射时，所需完整角直径为成人/儿童 \keyresult{\AdultInfinityAngle\degree}、小鸡 \keyresult{\ChickInfinityAngle\degree}。因此三种眼模型的核心设计量都是约 $20\degree$ 的覆盖角，而实体光源尺寸随工作距离变化。
  \item 外置镜片仍采用独立的 $10\D$ 参考工况：成人配合 $-10\D$ 镜片时需要 $17.004\D$ 调节，仍低于给定的 $18\D$ 上限；$-15\D$ 与 $-20\D$ 分别需要 $20.028\D$ 与 $22.787\D$，判为不可调节。
  \item 成人最大 $5\mm$ 瞳孔在 $10\D$ 离焦时产生 $0.835\mm$ 几何模糊斑。采用覆盖该离焦的保守光源后，目标区域捕获光线比例由 \FocusedCapture\% 降至 \DefocusedCapture\%，说明均匀覆盖与光利用率之间存在明确权衡。
  \item 五个聚焦 ZOS-API 案例中，独立模型与 OpticStudio 的视网膜边缘最大差异仅为 \keyresult{$\ZosMaxError\um$}；故意不调焦的成人 $10\D$ 案例得到 \ZosUnfocusedSpread\,mm 展宽，与解析预测一致。
\end{itemize}

这些结果支持第一阶段几何尺寸与离焦容差设计，但\textbf{不构成绝对视网膜辐照度或眼组织安全结论}。源文件未提供光源辐亮度、光谱带宽、曝光时间、组织透射率及完整解剖曲面；当前 ZOS 验证的是在真实 OpticStudio 中运行的近轴等效系统，而不是多曲面生理眼模型。

\clearpage
\tableofcontents
\clearpage
\pagenumbering{arabic}

\chapter{约 $20\degree$ 覆盖角决定系统的第一阶段几何设计}
\section{研究问题与决策目标}
本实验要回答的工程问题是：在小鸡、儿童与成人三个眼模型中，怎样选择 $650\,\mathrm{nm}$ 圆形面光源的实体尺寸或角尺寸，才能使其在调焦或给定离焦条件下覆盖目标后极部区域；同时，外置负镜片、瞳孔和眼轴变化会如何改变调节需求、模糊斑和相对照度分布。

报告把“覆盖”定义为\textbf{几何像盘覆盖指定后极部圆盘}。对于聚焦系统，像盘半径至少等于目标半径；对于离焦系统，同时给出几何最低覆盖角与均匀性保守覆盖角。该定义是工程筛选条件，并不等价于生物学上的“全视网膜剂量达标”。

\section{主要设计输出}
\begin{table}[H]
\centering
\caption{无外置镜片时的主要光源直径；括号内为给定调节上限下的可行性}
\label{tab:headline}
\small
\begin{tabularx}{\textwidth}{ccXXX}
\toprule
物方需求（D） & 物距（mm） & 小鸡直径（mm） & 儿童直径（mm） & 成人直径（mm） \\
\midrule
\HeadlineRows
\bottomrule
\end{tabularx}
\end{table}

\Cref{tab:headline} 显示，物方需求由 $60\D$ 增至 $120\D$ 时，物距由 $16.67\mm$ 缩短至 $8.33\mm$，所需实体光源直径恰好减半。由于无外置镜片时调节力等于物方需求，三个模型的全部请求工况均超过给定调节上限。表中的“超限”仅表示生理调节约束失败；几何光源尺寸仍可用于固定焦度或非生理模型的尺寸研究。

\begin{figure}[H]
\centering
\includegraphics[width=0.88\textwidth]{source_diameter_cn.png}
\caption{无外置镜片时，不同物方屈光需求下的光源直径。数据来自 \sourcefile{headline\_results.csv}。}
\label{fig:source-diameter}
\end{figure}

\Cref{fig:source-diameter} 应按七个有序离散工况阅读：$60\D$ 到 $120\D$ 对应 $16.67\mm$ 到 $8.33\mm$ 物距，步长为 $10\D$。三条曲线高度接近的原因不是三个眼完全相同，而是目标后极部直径与有效焦距近似按比例变化，使所需角覆盖接近。设计含义是：先以约 $20\degree$ 完整角作为视场约束，再根据安装距离换算实体光源尺寸。

\chapter{输入数据、参数定义与边界条件}
\section{源文件与证据范围}
唯一外部输入是用户提供的参数演示文稿，其原始副本保存在本实验的 \sourcefile{source} 目录，文件名为《小鸡和人眼光学仿真参数.pptx》。该文件的 SHA-256 为
\begin{center}
\small\ttfamily b8bcfe66c559081daecc29fa6426de1c297667848a2c1f7e60a5a3f198640e39。
\end{center}
模型参数被结构化到 \sourcefile{config/experiment.json}；所有报告数字由版本化 CSV/JSON 结果生成，而不是手工录入。源文件中的第5页“15 D”结合相邻的负镜片序列暂解释为 $-15\D$，这是需要实验人员确认的关键假设。

\section{三个等效眼模型}
\begin{table}[H]
\centering
\caption{模型参数与扫描范围}
\label{tab:eye-parameters}
\small
\begin{tabularx}{\textwidth}{lXXXXX}
\toprule
模型 & 有效焦距（mm） & 报告眼轴（mm） & 目标直径（mm） & 瞳孔扫描（mm） & 调节上限（D） \\
\midrule
30--45日龄小鸡 & 8.4 & 11.7 & 3.0 & 2.0, 2.5, 3.0, 3.5 & 15.69 \\
6岁儿童 & 16.7 & 23.0 & 6.0 & 2.0, 3.0, 4.0, 5.0 & 22.50 \\
18岁成人 & 16.7 & 23.6 & 6.0 & 2.0, 3.0, 4.0, 5.0 & 18.00 \\
\bottomrule
\end{tabularx}
\end{table}

基准眼屈光力定义为 $F_0=1000/f_0$（单位 D）。为了把源文件的完整眼轴与有效焦距之间的主平面偏置分离，模型把基准约化视网膜距离校准为 $t=1/F_0$。完整眼轴只用于轴长敏感性分析，不直接替换校准后的约化像距。

\section{实验变量、固定量与响应量}
\begin{table}[H]
\centering
\caption{实验设计中的变量分类}
\label{tab:variables}
\small
\begin{tabularx}{\textwidth}{p{2.6cm}p{5.3cm}X}
\toprule
类别 & 参数/水平 & 作用 \\
\midrule
固定量 & 波长 $650\,\mathrm{nm}$；均匀圆形朗伯面源；旋转对称近轴传播 & 限定本报告为单色几何光学和相对照度研究 \\
物方主扫描 & $60,70,80,90,100,110,120\D$；$s=1/|L|$ & 按用户要求以 $10\D$ 步长改变工作距离和调节需求 \\
外置镜片参考 & $10\D$ 物方需求；镜片 $0,-1,-3,-5,-10,-15,-20\D$ & 在物理元件顺序有效时测试负镜片与调节上限 \\
离焦扫描 & $-10,-5,0,2,5,10,15,20\D$ & 量化正负离焦的几何模糊 \\
个体参数 & 瞳孔直径、报告眼轴及其敏感性水平 & 量化通光口径与轴长偏差 \\
响应量 & 调节力、光源直径/面积、覆盖角、模糊斑、捕获比例、均匀性 & 支持几何尺寸和容差判断 \\
\bottomrule
\end{tabularx}
\end{table}

\chapter{等效眼、外置镜片与后极部的光学系统}
\section{示意光路说明模型部件及其顺序}
\begin{figure}[H]
\centering
\begin{tikzpicture}[x=1.12cm,y=1.05cm,>=Latex,font=\small]
  \draw[ReportGrey,thin] (-0.4,0)--(12.2,0);
  \draw[ReportOrange,very thick] (0,-1.5)--(0,1.5);
  \fill[ReportOrange] (0,1.5) circle(1.8pt) (0,-1.5) circle(1.8pt);
  \node[align=center,below=8pt] at (0,-1.5) {均匀圆形\\朗伯面源};

  % negative external lens
  \draw[ReportGold,fill=ReportWarm,thick]
    (3,-1.35) .. controls (3.45,-0.8) and (3.45,0.8) .. (3,1.35)
    -- (3.28,1.35) .. controls (2.83,0.8) and (2.83,-0.8) .. (3.28,-1.35) -- cycle;
  \node[align=center,below=8pt] at (3.14,-1.35) {外置负镜片\\$F_{\mathrm{ext}}$};

  % pupil
  \draw[ReportInk,very thick] (5,-1.35)--(5,-0.45) (5,0.45)--(5,1.35);
  \draw[ReportInk,<->] (4.78,-0.43)--(4.78,0.43);
  \node[below=8pt] at (5,-1.35) {瞳孔 $D_p$};

  % eye lens
  \draw[ReportBlue,fill=ReportLight,thick]
    (6.8,-1.45) .. controls (7.55,-0.75) and (7.55,0.75) .. (6.8,1.45)
    .. controls (6.05,0.75) and (6.05,-0.75) .. cycle;
  \node[align=center,below=8pt] at (6.8,-1.45) {可调等效眼\\$F_0+A$};

  % retina
  \draw[ReportPink,very thick] (10.45,-1.8) arc[start angle=-90,end angle=90,x radius=0.72,y radius=1.8];
  \draw[ReportPink,<->] (10.6,-1.5)--(10.6,1.5);
  \node[align=center,below=8pt] at (10.7,-1.8) {目标后极部\\直径 $D_r$};

  % rays
  \draw[ReportOrange,thick] (0,1.5)--(3,0.96)--(5,0.48)--(6.8,0.25)--(10.45,1.5);
  \draw[ReportBlue,thick] (0,0)--(3,0)--(5,0)--(6.8,0)--(10.45,0);
  \draw[ReportOrange,thick] (0,-1.5)--(3,-0.96)--(5,-0.48)--(6.8,-0.25)--(10.45,-1.5);

  \draw[<->,ReportGrey] (0,2.15)--(6.8,2.15) node[midway,fill=white] {物距 $s$};
  \draw[<->,ReportGrey] (3,1.75)--(6.8,1.75) node[midway,fill=white] {顶点距 $v$};
  \draw[<->,ReportGrey] (6.8,2.15)--(10.45,2.15) node[midway,fill=white] {约化视网膜距 $t$};
\end{tikzpicture}
\caption{带外置负镜片的等效眼照明系统示意图。图中光线仅用于说明部件、孔径和覆盖关系，不按真实曲率或角度比例绘制。}
\label{fig:lens-system}
\end{figure}

\Cref{fig:lens-system} 是本实验的建模边界。外置镜片位于光源与等效眼之间，其顶点距暂定为小鸡 $5\mm$、人眼 $12\mm$；瞳孔作为限束孔径；眼内全部解剖折射面折算为一个可调薄透镜；视网膜用固定约化距离处的目标圆盘表示。聚焦像盘尺寸在一阶模型中与瞳孔直径无关，但离焦模糊与通过量依赖瞳孔。

\section{该示意图不代表的内容}
图中未建立角膜前后表面、前房、晶状体梯度折射率、玻璃体、视网膜曲率、像差、散射和眼球转动。因此，光路图用于解释矩阵顺序和变量耦合，不能作为机械结构图或组织剂量图。后续多曲面模型应保持同一输入/输出定义，以便与本报告结果逐项比较。

\chapter{ABCD 模型把调焦条件转化为可审计的矩阵方程}
\section{光线状态与基本矩阵}
采用约化角光线状态
\begin{equation}
  \bm{r}=\begin{bmatrix}y\\u\end{bmatrix},\qquad u=n\theta,
\end{equation}
其中 $y$ 是光线高度，$n$ 是介质折射率，$\theta$ 为近轴角。距离 $d$、折射率 $n$ 中的平移矩阵和光焦度 $F$ 的薄透镜矩阵分别为
\begin{equation}
  \bm{T}(d,n)=
  \begin{bmatrix}1&d/n\\0&1\end{bmatrix},\qquad
  \bm{L}(F)=
  \begin{bmatrix}1&0\\-F&1\end{bmatrix}.
  \label{eq:basic-matrices}
\end{equation}

带外置镜片的眼前矩阵写为
\begin{equation}
  \bm{M}_{\mathrm{pre}}=
  \bm{T}(v,1)\,\bm{L}(F_{\mathrm{ext}})\,\bm{T}(s-v,1)
  =\begin{bmatrix}a&b\\c&d\end{bmatrix}.
\end{equation}
从光源到视网膜的总矩阵为
\begin{equation}
  \bm{M}_{\mathrm{tot}}=
  \bm{T}(t,1)\,\bm{L}(F_0+A)\,\bm{M}_{\mathrm{pre}}
  =\begin{bmatrix}A_m&B_m\\C_m&D_m\end{bmatrix}.
  \label{eq:total-matrix}
\end{equation}
这里 $A$ 表示附加调节力，避免与矩阵元素 $A_m$ 混淆。

\section{$B_m=0$ 同时给出调节力和横向放大率}
光源平面与视网膜共轭时，像高不依赖初始光线角，因此成像条件是
\begin{equation}
  B_m=0.
\end{equation}
代入\Cref{eq:total-matrix}，并利用 $t=1/F_0$，得到代码中使用的调节解
\begin{equation}
  A=\frac{d}{b},\qquad F_{\mathrm{eye}}=F_0+A.
  \label{eq:accommodation}
\end{equation}
在调焦状态，横向放大率为 $m=A_m$，覆盖目标直径 $D_r$ 所需的最小圆形光源直径和面积为
\begin{equation}
  D_s=\frac{D_r}{|m|},\qquad S_s=\pi\left(\frac{D_s}{2}\right)^2.
  \label{eq:source-size}
\end{equation}
数值质量控制要求所有扫描行满足 $|B_m|<10^{-12}\,\mathrm{m}$；21 行 $60$--$120\D$ 主扫描与 21 行 $10\D$ 外置镜片参考扫描全部通过。

\section{无穷远角覆盖与离焦模糊}
无外置镜片时，无穷远目标的完整角直径近似为
\begin{equation}
  \Theta_{\infty}=\frac{D_r}{t}.
\end{equation}
当眼存在离焦 $\Delta F$、瞳孔半径为 $r_p$ 时，一阶模糊半径为
\begin{equation}
  r_b=|t\,\Delta F|\,r_p.
  \label{eq:defocus-blur}
\end{equation}
因此仅追求“几何至少触及目标边界”的角半径和追求“源盘在离焦卷积后仍完全覆盖目标”的保守角半径分别为
\begin{equation}
  \theta_{\min}=\max\left(0,\frac{r_t-r_b}{t}\right),\qquad
  \theta_{\mathrm{con}}=\frac{r_t+r_b}{t}.
\end{equation}
后者提高均匀覆盖的稳健性，但会让更多光线落在目标圆盘之外。

\chapter{可复现实验设计同时覆盖几何、离焦、眼轴与照度}
\section{确定性工作流将每个结果绑定到配置和代码}
\begin{figure}[H]
\centering
\resizebox{0.97\textwidth}{!}{\begin{tikzpicture}[
  node distance=8mm and 11mm,
  box/.style={draw=ReportBlue,rounded corners=2pt,fill=ReportLight,minimum width=3.2cm,minimum height=0.9cm,align=center},
  zbox/.style={draw=ReportOrange,rounded corners=2pt,fill=ReportWarm,minimum width=3.2cm,minimum height=0.9cm,align=center},
  check/.style={draw=ReportPink,rounded corners=2pt,fill=ReportPink!8,minimum width=3.3cm,minimum height=0.9cm,align=center},
  >=Latex,font=\small]
  \node[box] (ppt) {原始 PPT\\参数与问题};
  \node[box,right=of ppt] (config) {结构化配置\\experiment.json};
  \node[box,right=of config] (model) {独立 ABCD 模型\\eye\_model.py};
  \node[box,below=of model] (sweep) {参数扫描与蒙特卡洛\\固定随机种子};
  \node[zbox,below=of config] (zos) {ZOS-API 自动建模\\OpticStudio 24.1};
  \node[check,below=of ppt] (tests) {23 项自动测试\\解析恒等式与边界};
  \node[check,below=14mm of sweep] (validate) {CSV/JSON 交叉验证\\误差与完整性检查};
  \node[box,left=of validate] (notebook) {执行型 Notebook\\图表与审计轨迹};
  \node[box,right=of validate] (report) {LaTeX/PDF 报告\\自检与 Git 提交};

  \draw[->] (ppt)--(config);
  \draw[->] (config)--(model);
  \draw[->] (model)--(sweep);
  \draw[->] (config)--(zos);
  \draw[->] (model)--(tests);
  \draw[->] (sweep)--(validate);
  \draw[->] (zos)--(validate);
  \draw[->] (tests)--(validate);
  \draw[->] (validate)--(notebook);
  \draw[->] (validate)--(report);
\end{tikzpicture}}
\caption{从原始参数到验证报告的可复现实验工作流。独立模型和 OpticStudio 路径在验证节点汇合。}
\label{fig:workflow}
\end{figure}

\Cref{fig:workflow} 中的关键原则是把独立计算与软件验证分开实现：Python 模型不依赖 ZOS-API，ZOS-API 验证器则用 C\# 创建系统并写出独立 CSV。两条路径只在 \sourcefile{validate\_results.py} 中比较，从而避免同一段计算逻辑同时生成“预测”和“验证”。

\section{主物方扫描与外置镜片参考扫描保持物理顺序有效}
主扫描包含 3 个眼模型和 7 个物方需求，共 $3\times7=21$ 行无外置镜片聚焦结果；外置镜片另在 $10\D$ 参考物距扫描 3 个眼模型和 7 个镜片水平，共 21 行；无穷远还有 $3\times7=21$ 行角覆盖结果。每行记录调节力、总眼屈光力、放大率、光源直径、光源面积、调节可行性与成像矩阵残差。

人眼外置镜片的暂定顶点距为 $12\mm$，而 $90$--$120\D$ 的物距只有 $11.11$--$8.33\mm$。若把这一区间与外置镜片直接组合，镜片将不再位于光源与眼之间，违反当前传递矩阵的元件顺序。因此报告不生成该无效组合，而保留几何有效的 $10\D$ 外置镜片参考扫描。

离焦扫描为 3 个眼模型、各自 4 个瞳孔、8 个离焦水平，共 96 行。眼轴敏感性为 3 个眼模型、各自 4 个瞳孔、3 个轴长水平，共 36 行。所有结果保留在 CSV 中，报告只抽取对工程判断最有用的切片。

\section{蒙特卡洛样本与均匀性指标}
面源点和瞳孔点均在圆盘内按面积均匀采样：若 $U\sim\mathcal{U}(0,1)$，则半径 $r=R\sqrt{U}$、方位角 $\phi=2\pi V$。每个案例使用 $600{,}000$ 条光线，随机种子分别为 20260814 和 20260815。

目标捕获比例定义为落入后极部目标圆盘的光线数除以发射总数；$p_{10}/\bar{N}$ 均匀性定义为目标圆盘内直方图非几何遮罩像素计数的第 10 百分位数除以均值。另记录达到均值一半以上的像素比例。指标描述相对分布，不含光源绝对功率。

\chapter{实体光源尺寸随物距变化，覆盖角保持近似稳定}
\section{60--120 D 扫描中实体光源直径随物距减半}
在无外置镜片条件下，成人和儿童从 $60\D$ 到 $120\D$ 分别需要 \AdultSixtyDDiameter\,mm 到 \AdultOneTwentyDDiameter\,mm 光源直径；小鸡需要 \ChickSixtyDDiameter\,mm 到 \ChickOneTwentyDDiameter\,mm。三种模型差异小于 $0.7\%$，说明第一阶段可以共享近似视场规格，但全部案例均超过给定调节上限。

\section{外置负镜片提高调节需求并略微增大光源尺寸}
\begin{figure}[H]
\centering
\includegraphics[width=0.9\textwidth]{adult_accommodation_cn.png}
\caption{成人模型在 $10\D$ 参考工况下，所需调节力随外置负镜片度数变化；虚线为 $18\D$ 调节上限。}
\label{fig:adult-heatmap}
\end{figure}

\Cref{fig:adult-heatmap} 表明，在固定 $10\D$ 参考物距下，负镜片持续推高成人调节需求；$-15\D$ 与 $-20\D$ 越过 $18\D$ 上限。该图不属于 $60$--$120\D$ 主扫描，而是为保留原有外置镜片敏感性证据设置的独立物理有效参考。

\begin{table}[H]
\centering
\caption{成人 $100\mm$ 物距（$10\D$）下的外置镜片扫描}
\label{tab:adult-lens}
\small
\begin{tabular}{rrrrl}
\toprule
镜片（D） & 所需调节（D） & 光源直径（mm） & 光源面积（mm$^2$） & 可行性 \\
\midrule
\AdultLensRows
\bottomrule
\end{tabular}
\end{table}

从 $0\D$ 到 $-10\D$，光源直径从 $35.928\mm$ 增至 $39.722\mm$，面积增至 $1239.2\,\mathrm{mm^2}$；同时调节力增至 $17.004\D$。因此外置镜片的影响不能只通过“度数相加”判断，还必须考虑镜片与眼之间的顶点距及其对横向放大率的影响。

\chapter{瞳孔不改变理想聚焦像尺寸，却主导离焦容差}
\section{离焦模糊随瞳孔直径和离焦绝对值近似线性增加}
\begin{figure}[H]
\centering
\includegraphics[width=\textwidth]{defocus_blur_cn.png}
\caption{三个眼模型的离焦—模糊斑关系；黑色虚线为各模型目标后极部直径。}
\label{fig:defocus}
\end{figure}

\Cref{fig:defocus} 与\Cref{eq:defocus-blur} 一致：对固定眼模型，模糊斑关于离焦符号对称，并与瞳孔半径成比例。聚焦状态下，像尺寸由放大率决定而与瞳孔无关；离焦后，瞳孔从通光量参数变为几何模糊的直接控制量。

\begin{table}[H]
\centering
\caption{$+10\D$ 离焦、各模型最大扫描瞳孔下的模糊与覆盖角}
\label{tab:defocus-summary}
\small
\begin{tabular}{lrrrr}
\toprule
模型 & 瞳孔（mm） & 模糊直径（mm） & 几何最小角（$\degree$） & 保守均匀角（$\degree$） \\
\midrule
\DefocusRows
\bottomrule
\end{tabular}
\end{table}

若只要求某些光线触及目标边界，离焦使“最小角”减小；若要求整个目标圆盘仍被源盘卷积覆盖，则应采用更大的“保守均匀角”。成人最大瞳孔在 $10\D$ 离焦时，二者分别为 $17.721\degree$ 与 $23.450\degree$。报告后续蒙特卡洛采用后者，因此它检验的是保守覆盖方案而不是最省光方案。

\section{眼轴偏差在小鸡模型中表现出更高的屈光敏感性}
物理眼轴变化 $\Delta l$ 先按像方折射率 $n'=1.336$ 换算为约化距离变化，再计算等效离焦：
\begin{equation}
  \Delta F_{\mathrm{eq}}=-\frac{\Delta l}{n'}F_0^2.
\end{equation}
\begin{table}[H]
\centering
\caption{各模型扫描范围内绝对等效离焦最大的眼轴案例；采用最大扫描瞳孔}
\label{tab:axial}
\small
\begin{tabular}{lrrrr}
\toprule
模型 & 瞳孔（mm） & 眼轴偏差（mm） & 等效离焦（D） & 模糊直径（mm） \\
\midrule
\AxialRows
\bottomrule
\end{tabular}
\end{table}

小鸡较短的有效焦距对应更高的基准屈光力，因此同量级轴长偏差产生更大的等效离焦。\Cref{tab:axial} 中小鸡 $-1.2\mm$ 轴长偏差对应 $+12.730\D$，最大瞳孔模糊直径 $0.374\mm$。这提示小鸡系统的轴长实测与个体分层比成人模型更加重要。

\chapter{保守离焦光源改善覆盖稳健性，但降低目标捕获效率}
\section{聚焦案例在目标圆盘内保持近似均匀分布}
\begin{figure}[H]
\centering
\includegraphics[width=\textwidth]{monte_carlo_cn.png}
\caption{成人眼两个蒙特卡洛案例的视网膜相对光线计数。白色圆为直径 $6\mm$ 的目标后极部区域。}
\label{fig:monte-carlo}
\end{figure}

\Cref{fig:monte-carlo} 左图使用 $100\mm$ 物距下的最小聚焦光源，捕获比例为 \FocusedCapture\%，$p_{10}/\bar{N}=\FocusedUniformity$。随机采样造成像素级噪声，但没有系统性中心暗斑或边缘缺口。

右图使用覆盖 $+10\D$ 离焦的保守角尺寸，$p_{10}/\bar{N}=\DefocusedUniformity$，目标圆盘仍被照亮；但只有 \DefocusedCapture\% 的发射光线落入目标区域，其余光线落到目标圆盘之外。因未建模全视网膜边界与组织吸收，这部分光线不能被简单称为“损失”，也不能据此评价安全性。

\section{蒙特卡洛结果的统计解释边界}
本实验的随机性只来自几何采样，模型参数本身不是随机变量。固定种子使图和汇总指标可重复，但没有为生理个体差异建立概率分布。因此这些统计量是给定模型和给定光源的数值积分近似，不是对小鸡或人群的置信区间。

\chapter{OpticStudio 24.1 在六个系统上复现了解析边缘与离焦展宽}
\section{ZOS-API 验证是独立软件路径，而不是 Python 结果回显}
C\# 验证器通过 \sourcefile{ZOSAPI\_Initializer} 连接已安装的 OpticStudio，确认 API 许可证有效后，为每个案例新建顺序模式系统。系统使用 Paraxial 表面表示外置镜片和等效眼透镜，保存为独立 \sourcefile{.zos} 文件，并以标准化批量真实光线追迹接口读取像面光线高度。

验证共含五个调焦案例和一个故意不调焦案例。每个案例使用五条具有代表性的边缘/孔径光线进行确定性几何闭合检查。该样本足以验证一阶矩阵边缘位置和离焦展宽，但不等同于高密度光斑、像差或杂散光分析。

\begin{figure}[H]
\centering
\includegraphics[width=0.9\textwidth]{zos_validation_cn.png}
\caption{五个聚焦案例的独立目标边缘与 OpticStudio 光线追迹像高。圆点和叉号重合。}
\label{fig:zos-validation}
\end{figure}

\Cref{fig:zos-validation} 中的重合不是视觉分辨率不足造成的模糊判断；\Cref{tab:zos} 给出数值误差。聚焦边缘误差和 RMS 展宽均处于双精度浮点舍入量级。

\begin{table}[H]
\centering
\caption{OpticStudio ZOS-API 六个验证案例的数值结果}
\label{tab:zos}
\scriptsize
\begin{tabularx}{\textwidth}{p{3.1cm}c r r r r r}
\toprule
案例 & 调焦 & 目标边缘（mm） & ZOS边缘（mm） & 边缘误差（$\um$） & RMS（$\um$） & 展宽（mm） \\
\midrule
\ZosRows
\bottomrule
\end{tabularx}
\end{table}

故意不调焦的成人 $10\D$ 案例保持基准眼屈光力，解析预测像高总展宽为 $0.82665\mm$，OpticStudio 观测值为 \ZosUnfocusedSpread\,mm。这个负对照证明验证器能够检测非零离焦，而不是无条件输出目标边缘。

\section{验证结论只覆盖当前模型阶次}
交叉验证说明 Python ABCD 实现、外置镜片顺序、调节解和 ZOS Paraxial 系统在一阶光学上相互一致。它不证明真实眼的角膜/晶状体面形、像差、散射或组织辐照度正确。将来切换到多曲面模型时，应增加光斑 RMS、波前误差、畸变、相对照度、非序列探测器和公差统计等新的验证量。

\chapter{限制、不确定性与稳健性检查决定结果的使用边界}
\section{已完成的稳健性检查}
\begin{enumerate}[leftmargin=2.5em]
  \item \textbf{解析恒等式测试：}无外置镜片时，调节力等于物方屈光需求；零离焦时模糊为零；标称眼轴时轴长模糊为零。
  \item \textbf{矩阵残差检查：}21 行 $60$--$120\D$ 主扫描和 21 行外置镜片参考结果的成像矩阵 $B_m$ 残差全部低于 $10^{-12}\,\mathrm{m}$。
  \item \textbf{符号和边界检查：}负镜片应提高调节需求；调节可行性要求 $0\le A\le A_{\max}$；顶点距必须位于光源与眼之间。
  \item \textbf{负对照：}成人 $10\D$ 故意不调焦案例在独立模型和 ZOS 中同时给出 $0.82665\mm$ 展宽。
  \item \textbf{确定性复现：}蒙特卡洛使用固定种子；配置和源 PPT 具有 SHA-256；实验记录通过不可覆盖脚本写入 Git。
\end{enumerate}

\section{可能改变设计结论的缺失证据}
\begin{longtable}{p{3.0cm}p{5.0cm}p{5.7cm}}
\caption{当前不确定性及其潜在影响}\label{tab:limitations}\\
\toprule
缺失或假设 & 当前处理 & 对结论的可能影响 \\
\midrule
\endfirsthead
\toprule
缺失或假设 & 当前处理 & 对结论的可能影响 \\
\midrule
\endhead
绝对光源功率与曝光 & 未提供，报告只输出相对光线计数 & 无法计算视网膜辐照度、热剂量、光化学剂量或安全裕量 \\
光谱带宽与组织透射 & 固定单色 $650\,\mathrm{nm}$，未建组织吸收 & 不能外推到宽谱光源或比较组织层吸收 \\
真实解剖曲面和折射率 & 用一个等效薄透镜代替 & 可能改变像差、场曲、畸变、边缘照度和有效焦距 \\
外置镜片顶点距 & 小鸡 $5\mm$、人眼 $12\mm$ 暂定 & 会改变调节力与放大率，尤其在高负度数下 \\
超近距元件顺序 & $90$--$120\D$ 人眼物距短于暂定 $12\mm$ 顶点距 & 不生成外置镜片组合；需重定义机械位置后才能建模 \\
第5页 15 D 符号 & 按序列解释为 $-15\D$ & 若原意为 $+15\D$，相关调节结论将反向变化 \\
个体差异分布 & 只做离散眼轴和瞳孔扫描 & 不能给出人群覆盖率或统计置信区间 \\
眼球转动和全视网膜包络 & 目标固定为后极部圆盘 & “全视网膜”实际机械视场可能大于当前约 $20\degree$ 后极部覆盖角 \\
散射、遮挡与杂散光 & 未进入非序列模型 & 无法评估非成像路径、虹膜/眼睑遮挡和离轴安全风险 \\
\bottomrule
\end{longtable}

\section{结论应被表述为描述性和模型条件性结果}
本报告的数值是“在给定参数和一阶模型下”的描述性预测，并通过同阶 OpticStudio 系统交叉验证。报告不包含因果推断、临床有效性推断或生物安全认证。任何硬件定型之前，应把实测参数、完整几何和辐射度约束纳入同一可复现管线。

\chapter{一条命令可重建结果、图、Notebook、Zemax 系统与 PDF}
\section{软件环境}
已验证环境包括 Windows、Python 3.11、NumPy 2.4.6、pandas 3.0.5、Matplotlib 3.11.1、Ansys Zemax OpticStudio 24.1、有效 ZOS-API 许可证、MiKTeX 25.12 和 XeLaTeX。中文图表和 PDF 正文使用 Windows 宋体；PDF 构建脚本检查宋体字体是否嵌入。

\section{完整实验复现}
在仓库根目录安装依赖：
\begin{lstlisting}
python -m pip install -e ".[dev,analysis]"
\end{lstlisting}
然后运行主实验：
\begin{lstlisting}
powershell -ExecutionPolicy Bypass `
  -File experiments\eye_illumination\run_all.ps1 `
  -PythonPath "C:\path\to\python.exe" `
  -OpticStudioDir "C:\Program Files\Ansys Zemax OpticStudio 2024 R1.00"
\end{lstlisting}
该命令依次重建独立模型、CSV/JSON、五张基础图、C\# ZOS-API 验证器、六个 Zemax 系统、验证报告、Notebook、23 项测试和 HTML 报告。

\section{LaTeX 报告复现}
完成主实验后运行：
\begin{lstlisting}
powershell -ExecutionPolicy Bypass `
  -File experiments\eye_illumination\report\latex\build_report.ps1 `
  -PythonPath "C:\path\to\python.exe"
\end{lstlisting}
脚本首先从 CSV/JSON 生成 \sourcefile{generated\_results.tex}，再用宋体重绘五张中文图；随后 XeLaTeX 连续编译三次以稳定目录和交叉引用，最后检查 PDF 页数、A4 页面、字体、图片实例、必要章节文本和稀疏页，并生成联系表预览与机器可读 QA JSON。

\section{版本化输出与数据血缘}
\begin{longtable}{p{6.4cm}p{8.5cm}}
\caption{主要文件及其用途}\label{tab:artifacts}\\
\toprule
相对路径 & 用途 \\
\midrule
\endfirsthead
\toprule
相对路径 & 用途 \\
\midrule
\endhead
\sourcefile{config/experiment.json} & 参数、扫描水平、假设与随机种子的唯一配置源 \\
\sourcefile{eye\_model.py} & ABCD、调节、离焦、眼轴与蒙特卡洛核心 \\
\sourcefile{run\_experiment.py} & 全部扫参、CSV/JSON 和基础图生成 \\
{\footnotesize\nolinkurl{zemax/ZosApiEyeValidation.cs}} & 真实 OpticStudio 系统创建与批量光线追迹 \\
\sourcefile{results/} & 21 行主聚焦、21 行外置镜片参考、96 行离焦、36 行眼轴、21 行无穷远及验证结果 \\
{\footnotesize\nolinkurl{notebooks/eye_illumination_analysis.ipynb}} & 已执行分析 Notebook 和审计轨迹 \\
\sourcefile{report/latex/} & 本报告源码、数据生成器、中文图生成器与 PDF 自检 \\
{\footnotesize\nolinkurl{experiments/runs/eye-illumination-60-120d-v2.json}} & 本次 $60$--$120\D$ 更新的不可覆盖里程碑记录 \\
\bottomrule
\end{longtable}

公开仓库地址为 \url{https://github.com/dongzhaohe321418-lab/zemax-mcp}。大体积 PPT、PNG、Zemax 与 PDF 文件由 Git LFS 管理，文本代码和 CSV/JSON 由普通 Git 历史管理。

\chapter{下一阶段应从实测辐射度和多曲面眼模型开始}
\section{证据支持的优先工作}
\begin{enumerate}[leftmargin=2.5em]
  \item \textbf{补齐绝对辐射度：}测量光源光谱辐亮度、曝光时间、工作距离和调制方式，加入角膜到视网膜的光谱透射，输出辐照度与剂量，而不再只输出相对计数。
  \item \textbf{升级解剖模型：}用角膜前后表面、房水、晶状体、玻璃体和曲面视网膜替换等效薄透镜；保留当前一阶模型作为回归基线。
  \item \textbf{确认机械几何：}实测外置镜片顶点距、倾斜、偏心、光源尺寸及眼球位置分布，并建立制造/装配公差蒙特卡洛。
  \item \textbf{扩展全视网膜包络：}将目标从固定后极部圆盘扩展到眼球转动、视轴偏差和边缘视网膜共同形成的三维包络。
  \item \textbf{进入非序列验证：}加入虹膜、眼睑、外壳、散射和杂散光路径，使用探测器面验证均匀性与非目标区域暴露。
  \item \textbf{外部实验验证：}用模型眼或实验相机测量照度图，对放大率、模糊斑、边缘照度和捕获效率逐项校准。
\end{enumerate}

\section{仍需实验负责人确认的问题}
\begin{itemize}[leftmargin=2.2em]
  \item 第5页“15 D”是否确为 $-15\D$？
  \item 小鸡和人眼外置镜片的实际顶点距、倾斜和有效口径是多少？
  \item “全视网膜覆盖”的验收标准是几何覆盖、最低照度、均匀性，还是累积安全剂量？
  \item 是否有每只眼的眼轴、屈光状态、瞳孔和调节力原始样本数据，可用于概率模型？
  \item 光源是否为裸朗伯面源，还是包含扩散片、导光结构或准直透镜？
\end{itemize}

\chapter*{结论}
\addcontentsline{toc}{chapter}{结论}
在当前参数和近轴等效眼假设下，三个模型覆盖指定后极部的完整角需求均约为 $20\degree$。按用户要求，主扫描现为 $60$--$120\D$、步长 $10\D$：成人与儿童的所需圆形光源直径从约 $5.99\mm$ 降至 $2.99\mm$，小鸡从约 $5.95\mm$ 降至 $2.98\mm$；但全部 21 个请求案例均超过给定调节上限。负外置镜片仅在独立 $10\D$ 参考工况中评估，$-15\D$ 及更负镜片越过成人调节上限。瞳孔主要影响离焦模糊而非理想聚焦像尺寸；保守离焦设计可维持覆盖，但会降低目标区域捕获效率。

独立 Python 模型、自动测试、确定性蒙特卡洛与真实 OpticStudio 24.1 ZOS-API 系统共同形成了可复现证据链。六个 ZOS 案例验证了一阶边缘和离焦展宽，但下一阶段必须补齐绝对辐射度、多曲面解剖模型、机械公差、眼球转动与非序列杂散光，才能把本报告从几何预设计推进到硬件和生物安全决策。

\appendix
\chapter{结果字段定义}
\begin{longtable}{p{5.2cm}p{9.7cm}}
\caption{主要结果字段及单位}\label{tab:fields}\\
\toprule
字段 & 定义 \\
\midrule
\endfirsthead
\toprule
字段 & 定义 \\
\midrule
\endhead
\sourcefile{source\_demand\_D} & 物方屈光需求绝对值；有限物距 $s=1/L$ \\
\sourcefile{accommodation\_D} & 为满足 $B_m=0$ 所需的眼附加屈光力 \\
\sourcefile{source\_diameter\_mm} & 聚焦像盘覆盖目标圆盘所需的最小实体光源直径 \\
\sourcefile{feasible\_accommodation} & $0\le A\le A_{\max}$ 时为真 \\
\sourcefile{imaging\_B\_residual\_m} & 调焦后总系统矩阵的 $B_m$ 元素 \\
\sourcefile{blur\_diameter\_mm} & 边缘瞳孔光线形成的几何离焦模糊直径 \\
{\footnotesize\nolinkurl{uniform_conservative_angular_diameter_deg}} & 离焦卷积后仍完整覆盖目标圆盘所需的保守完整角 \\
\sourcefile{captured\_ray\_fraction} & 落入目标后极部圆盘的光线数除以发射光线总数 \\
\sourcefile{p10\_to\_mean\_uniformity} & 目标圆盘像素计数第10百分位数与均值之比 \\
\sourcefile{rms\_spread\_um} & ZOS 验证光线像高相对均值的 RMS 展宽 \\
\bottomrule
\end{longtable}

\chapter{PDF 自检清单}
构建脚本在报告交付前执行以下机器检查：
\begin{itemize}[leftmargin=2.2em]
  \item PDF 至少 15 页，所有页面宽高比符合 A4；
  \item 至少包含五个栅格图像实例，且能提取到技术摘要、ABCD、OpticStudio、蒙特卡洛、可复现工作流、限制和附录等关键文本；
  \item PDF 字体清单中存在 SimSun/宋体子集；
  \item XeLaTeX 日志中不存在未解析引用或 Overfull \verb|\hbox|；
  \item 生成首页、四分之一页、中间页、四分之三页和末页的联系表，并由人工检查标题、表格、图片、示意图和页脚；
  \item 重新运行 23 项 Python 测试，检查 Git 差异和敏感信息，再提交并推送 GitHub。
\end{itemize}

\begin{thebibliography}{9}
\bibitem{ppt} 用户提供文稿：《小鸡和人眼光学仿真参数》，原始 PPTX，2026。
\bibitem{config} 本项目：\sourcefile{experiments/eye\_illumination/config/experiment.json}，结构化实验参数与假设。
\bibitem{model} 本项目：\sourcefile{experiments/eye\_illumination/eye\_model.py}，ABCD 等效眼与蒙特卡洛模型。
\bibitem{zos} 本项目：\sourcefile{results/zemax/zosapi\_validation.csv} 与六个 \sourcefile{.zos} 系统，OpticStudio 24.1 ZOS-API 验证。
\bibitem{notebook} 本项目：\sourcefile{notebooks/eye\_illumination\_analysis.ipynb}，已执行分析 Notebook。
\end{thebibliography}

\end{document}
